% ============ 3. Objetivos (D3) ============
\section{Objetivos}

\subsection{Objetivo general}

\textbf{Cuantificar y auditar de forma reproducible la fiabilidad,
comparabilidad y equidad del padrón de investigadores reconocidos de
MinCiencias (convocatorias 2013--2021), a partir de los datos abiertos
oficiales y su cruce con la producción de grupos, y entregar un marco de
indicadores y un tablero de auditoría que orienten la mejora del observatorio
de capital humano científico en Colombia.}

El objetivo general responde exactamente al problema planteado: no es más
ancho (se limita al padrón y su cruce con producción, 2013--2021) ni más
angosto (abarca las tres dimensiones de la pregunta: fiabilidad,
comparabilidad y equidad).

\subsection{Objetivos específicos}

Los tres objetivos específicos son secuenciales: cada uno habilita al
siguiente, y cumplidos todos agotan el objetivo general sin solaparse.

\begin{enumerate}
  \item \textbf{Cuantificar la calidad del registro} del padrón de
        investigadores reconocidos de MinCiencias (2013--2021) según las
        dimensiones de completitud, unicidad, validez y consistencia, usando
        el dataset abierto oficial \texttt{bqtm-4y2h} y su artefacto
        consolidado.
        \emph{Evidencia de cumplimiento:} tabla de indicadores de calidad por
        convocatoria con conteos verificables (filas reales vs.\ declaradas,
        duplicados de \texttt{id\_persona}, atípicos de edad, errores de
        codificación) y código reproducible que la genera.

  \item \textbf{Estimar la dinámica longitudinal} (retención, transición de
        categoría y tiempo hasta la salida) de los investigadores entre
        convocatorias consecutivas, contrastando cadenas de Markov y modelos
        de supervivencia.
        \emph{Evidencia de cumplimiento:} matrices de transición con IC,
        curvas de Kaplan--Meier y un modelo de Cox con covariables (área,
        género, región), todos reproducibles sobre los datos abiertos.

  \item \textbf{Cuantificar la equidad} del sistema en sus dimensiones
        territorial y de género, y de diversidad (etnia, discapacidad),
        con indicadores de concentración (HHI/Gini/Theil) y razones de
        representación frente al censo DANE 2018; la condición de víctima del
        conflicto, que el censo no captura, se mide como prevalencia interna
        del padrón en la ventana disponible (2021).
        \emph{Evidencia de cumplimiento:} tabla de indicadores de equidad con
        intervalos de confianza y figura resumen por convocatoria.
\end{enumerate}

Cada objetivo especifica \emph{con qué evidencia se declarará cumplido}, y la
secuencia es visible: sin el registro limpio (OE1) no hay panel longitudinal
fiable (OE2), y sin ambos no hay lectura de equidad comparable (OE3). El
marco de indicadores y el tablero de auditoría se derivan de los tres.
